\chapter{Error Handling}
\label{ch:error-handling}

\begin{quote}
\textit{``Errors should never pass silently. Unless explicitly silenced.''} --- The Zen of Python
\end{quote}

Novus treats errors as values, not exceptions. Every operation that can fail returns a \texttt{Result} type, making error handling explicit and impossible to ignore. This chapter covers Novus's error handling philosophy, the standard error types, and patterns for writing robust code.

\section{Philosophy: Errors as Values}
\label{sec:error-philosophy}

Unlike languages with exceptions, Novus requires explicit handling of errors:

\begin{itemize}
    \item \textbf{No hidden control flow} --- You always know when a function can fail
    \item \textbf{No uncaught exceptions} --- Errors can't silently propagate up the stack
    \item \textbf{Compile-time enforcement} --- Ignoring a \texttt{Result} produces a warning
    \item \textbf{Zero runtime cost} --- No exception tables or unwinding machinery
\end{itemize}

\begin{comingfromc}
In C, errors are typically indicated by return values (\texttt{-1}, \texttt{NULL}) or global \texttt{errno}. This leads to bugs when callers forget to check.

In Novus, \texttt{Result<T, E>} forces you to handle both success and failure cases. You cannot accidentally use an error value as if it were successful.

\begin{tabular}{ll}
\textbf{C} & \textbf{Novus} \\
\texttt{if (ptr == NULL)} & \texttt{match result \{ ... \}} \\
\texttt{if (fd < 0)} & \texttt{let file = open(path)?} \\
\texttt{errno} & \texttt{Result::Err(DosError::NotFound)} \\
\end{tabular}
\end{comingfromc}

\section{The Result Type}
\label{sec:result-type}

The core error handling type is \texttt{Result<T, E>}:

\begin{lstlisting}
enum Result<T, E> {
    Ok(T),    // Success with value
    Err(E),   // Failure with error
}
\end{lstlisting}

Every function that can fail returns a \texttt{Result}:

\begin{lstlisting}
fn parse_number(s: *u8) -> Result<i32, ParseError> {
    // Returns Ok(value) or Err(error)
}

fn open_file(path: *u8) -> Result<FileHandle, DosError> {
    // Returns Ok(handle) or Err(error)
}
\end{lstlisting}

\section{Working with Results}
\label{sec:working-with-results}

\subsection{Pattern Matching}

The most explicit way to handle results:

\begin{lstlisting}
let result = open_file("data.txt")

match result {
    Result::Ok(file) => {
        // Use the file
        let data = read_all(&file)
        close(file)
    },
    Result::Err(err) => {
        println("Failed to open file")
        println(err.message())
    }
}
\end{lstlisting}

\subsection{The ? Operator}

For functions that return \texttt{Result}, the \texttt{?} operator provides concise error propagation:

\begin{lstlisting}
fn process_file(path: *u8) -> Result<String, DosError> {
    let file = open_file(path)?       // Returns Err early if failed
    let data = read_all(&file)?       // Returns Err early if failed
    let parsed = parse_data(&data)?   // Returns Err early if failed
    Result::Ok(parsed)
}
\end{lstlisting}

The \texttt{?} operator:
\begin{enumerate}
    \item Checks if the result is \texttt{Err}
    \item If \texttt{Err}, returns immediately with that error
    \item If \texttt{Ok}, unwraps the value and continues
\end{enumerate}

\subsection{Providing Defaults}

Use \texttt{unwrap\_or} for safe fallback values:

\begin{lstlisting}
// ... (illustrative - with load_config defined elsewhere)
let config = load_config("app.cfg")
let value = config.unwrap_or(default_config())

// Or with a closure for lazy evaluation (note space between pipes)
let value = config.unwrap_or_else(| | compute_default())
\end{lstlisting}

\subsection{Panicking on Error}

Use \texttt{unwrap()} only when failure is impossible or indicates a bug:

\begin{lstlisting}
// Only use when you KNOW it can't fail
let value = result.unwrap()  // Panics on Err!

// Better: use expect() with a message
let value = result.expect("Config file must exist")
\end{lstlisting}

\textbf{Warning:} Panics crash the program. Use sparingly.

\section{Standard Error Types}
\label{sec:standard-errors}

Novus provides error types for each AmigaOS subsystem.

\subsection{DosError}

Errors from the DOS library (file operations, paths):

\begin{lstlisting}
pub enum DosError {
    NotFound,           // File/directory not found
    AlreadyExists,      // Object already exists
    DirNotFound,        // Directory not found
    DirectoryNotEmpty,  // Can't delete non-empty dir
    PermissionDenied,   // Access denied
    WriteProtected,     // File is write-protected
    ReadProtected,      // File is read-protected
    DiskFull,           // No space on disk
    NoDisk,             // No disk in drive
    DeviceNotMounted,   // Device not available
    InvalidLock,        // Bad file lock
    Unknown,            // Unmapped error code
}
\end{lstlisting}

\begin{lstlisting}
from std::error::errors import DosError, dos_last_error

fn read_config() -> Result<String, DosError> {
    let file = match open("config.txt", MODE_OLDFILE) {
        Result::Ok(f) => f,
        Result::Err(_) => return Result::Err(dos_last_error())
    }
    // ...
}
\end{lstlisting}

\subsection{ExecError}

Errors from the Exec library (memory, tasks, signals):

\begin{lstlisting}
pub enum ExecError {
    NoMem,              // Memory allocation failed
    BadSignal,          // Invalid signal number
    NullReturn,         // Unexpected null return
    BadTask,            // Invalid task pointer
    BadPort,            // Invalid message port
    SignalInUse,        // Signal already allocated
    TaskNotFound,       // FindTask failed
    LibraryNotFound,    // OpenLibrary failed
    IndexOutOfBounds,   // Array bounds violation
    CapacityExceeded,   // Fixed buffer full
}
\end{lstlisting}

\begin{lstlisting}
from std::error::errors import ExecError

fn allocate_buffer(size: u32) -> Result<*u8, ExecError> {
    let ptr = AllocMem(size, MEMF_PUBLIC)
    if ptr == null {
        return Result::Err(ExecError::NoMem)
    }
    Result::Ok(ptr)
}
\end{lstlisting}

\subsection{IntuitionError}

Errors from the Intuition library (windows, screens, gadgets):

\begin{lstlisting}
pub enum IntuitionError {
    WindowOpenFailed,     // OpenWindow failed
    ScreenOpenFailed,     // OpenScreen failed
    GadgetCreateFailed,   // Gadget creation failed
    MenuCreateFailed,     // Menu creation failed
    InvalidWindow,        // Bad window pointer
    InvalidScreen,        // Bad screen pointer
    NoIDCMP,              // IDCMP setup failed
    BadDisplayMode,       // Invalid display mode
}
\end{lstlisting}

\subsection{GraphicsError}

Errors from the Graphics library (bitmaps, sprites, fonts):

\begin{lstlisting}
pub enum GraphicsError {
    BitMapAllocFailed,    // BitMap allocation failed
    FontOpenFailed,       // OpenFont failed
    InvalidBitMap,        // Bad BitMap pointer
    InvalidRastPort,      // Bad RastPort pointer
    InvalidSpriteChannel, // Bad sprite channel (0-7)
    SpriteChannelInUse,   // Channel already allocated
    NoSpritesAvailable,   // All sprites in use
}
\end{lstlisting}

\section{The NovusError Wrapper}
\label{sec:novus-error}

When a function may fail with errors from multiple subsystems, use \texttt{NovusError}:

\begin{lstlisting}
pub enum NovusError {
    Dos(DosError),
    Exec(ExecError),
    Intuition(IntuitionError),
    Graphics(GraphicsError),
}
\end{lstlisting}

The \texttt{From} trait implementations enable automatic conversion with the \texttt{?} operator:

\begin{lstlisting}
fn load_and_display(path: *u8) -> Result<(), NovusError> {
    let data = read_file(path)?       // DosError -> NovusError
    let screen = open_screen()?       // IntuitionError -> NovusError
    let bitmap = load_bitmap(&data)?  // GraphicsError -> NovusError
    display(screen, bitmap)
    Result::Ok(())
}
\end{lstlisting}

\section{Error Messages}
\label{sec:error-messages}

All standard error types implement a \texttt{message()} method:

\begin{lstlisting}
match result {
    Result::Err(err) => {
        println("Error: ")
        println(err.message())
    },
    _ => {}
}
\end{lstlisting}

Error messages are human-readable strings:
\begin{itemize}
    \item \texttt{DosError::NotFound} $\rightarrow$ \texttt{"Object not found"}
    \item \texttt{ExecError::NoMem} $\rightarrow$ \texttt{"Memory allocation failed"}
    \item \texttt{IntuitionError::WindowOpenFailed} $\rightarrow$ \texttt{"Failed to open window"}
\end{itemize}

\section{Error Conversion}
\label{sec:error-conversion}

Convert between error types using provided functions:

\begin{lstlisting}
from std::error::errors import novus_error_from_dos

fn wrapper() -> Result<(), NovusError> {
    let file = match open("file.txt", MODE_OLDFILE) {
        Result::Ok(f) => f,
        Result::Err(e) => return Result::Err(novus_error_from_dos(e))
    }
    Result::Ok(())
}
\end{lstlisting}

Or get raw error codes for debugging:

\begin{lstlisting}
from std::error::errors import dos_error_to_code

let code = dos_error_to_code(DosError::NotFound)  // Returns 205
\end{lstlisting}

\section{Patterns and Best Practices}
\label{sec:error-patterns}

\subsection{Early Return Pattern}

Validate inputs and return early on error:

\begin{lstlisting}
fn process(data: *u8, len: u32) -> Result<(), Error> {
    if data == null {
        return Result::Err(Error::InvalidInput)
    }
    if len == 0 {
        return Result::Err(Error::EmptyData)
    }

    // Main logic here
    Result::Ok(())
}
\end{lstlisting}

\subsection{Error Context}

Add context when propagating errors:

\begin{lstlisting}
// ... (illustrative - error context example)
fn load_user_config() -> Result<Config, ConfigError> {
    let path = get_config_path()?
    let file = open(path, MODE_OLDFILE)
        .map_err(|e: DosError| ConfigError::FileNotFound)?
    let data = read_all(&file)
        .map_err(|e: DosError| ConfigError::ReadFailed)?
    parse_config(&data)
}
\end{lstlisting}

\subsection{Fallback Chains}

Try multiple approaches, falling back on failure:

\begin{lstlisting}
// ... (illustrative - fallback chains)
fn find_config() -> Result<Config, ConfigError> {
    // Try user config first
    match load_from("ENV:app.cfg") {
        Result::Ok(cfg) => return Result::Ok(cfg),
        Result::Err(_) => {},
    }

    // Fall back to system config
    match load_from("SYS:Prefs/app.cfg") {
        Result::Ok(cfg) => return Result::Ok(cfg),
        Result::Err(_) => {},
    }

    // Use defaults
    Result::Ok(default_config())
}
\end{lstlisting}

\subsection{Cleanup on Error}

Use \texttt{defer} to ensure cleanup even when returning early:

\begin{lstlisting}
fn process_file(path: *u8) -> Result<Data, DosError> {
    let file = open(path, MODE_OLDFILE)?
    defer { close(file) }  // Always runs

    let header = read_header(&file)?  // May return early
    let body = read_body(&file)?      // May return early

    Result::Ok(combine(header, body))
}
\end{lstlisting}

\section{Creating Custom Errors}
\label{sec:custom-errors}

Define domain-specific error types:

\begin{lstlisting}
pub enum ImageError {
    InvalidFormat,
    UnsupportedDepth(u8),
    CorruptedData,
    TooLarge(u32, u32),
}

impl ImageError {
    pub fn message(&self) -> *u8 {
        match *self {
            ImageError::InvalidFormat => return "Invalid image format",
            ImageError::UnsupportedDepth(_) => return "Unsupported color depth",
            ImageError::CorruptedData => return "Image data corrupted",
            ImageError::TooLarge(_, _) => return "Image dimensions too large",
        }
    }
}
\end{lstlisting}

\section{Summary}
\label{sec:error-summary}

\begin{itemize}
    \item \textbf{Use \texttt{Result<T, E>}} for all fallible operations
    \item \textbf{Pattern match} for explicit handling of all cases
    \item \textbf{Use \texttt{?}} to propagate errors concisely
    \item \textbf{Use \texttt{unwrap\_or}} for safe defaults
    \item \textbf{Avoid \texttt{unwrap()}} except when failure is impossible
    \item \textbf{Use standard errors} (DosError, ExecError, etc.) for AmigaOS operations
    \item \textbf{Use NovusError} when crossing subsystem boundaries
    \item \textbf{Use \texttt{defer}} to ensure cleanup on error paths
\end{itemize}

Error handling in Novus is explicit but ergonomic. The \texttt{?} operator and \texttt{Result} combinators make it easy to write code that correctly handles failures without obscuring the happy path.
